\section{Related Work}
\label{sec:related_work}

Our work occupies a unique intersection between parameter-efficient adaptation, weight merging, dynamic routing, and non-linear biological systems.

\subsection{Parameter-Efficient Fine-Tuning and Model Merging}
Parameter-Efficient Fine-Tuning (PEFT) has revolutionized model adaptation, enabling the modification of massive pre-trained models using a tiny fraction of trainable parameters. Notable PEFT approaches include Low-Rank Adaptation (LoRA) \cite{hu2021lora}, Prefix-Tuning \cite{li2021prefix}, Prompt Tuning \cite{lester2021power}, and BitFit \cite{zaken2021bitfit}. Among these, LoRA has emerged as the de facto standard due to its zero-latency inference after merging and modular design.

With the proliferation of task-specific adapters, researchers have proposed static model merging to combine multiple capabilities into a single set of weights without retraining. 
Model Soups \cite{wortsman2022model} average weights of models fine-tuned on the same task. For disparate tasks, Task Arithmetic \cite{ilharco2022editing} defines task vectors that can be linearly added or subtracted. 
However, simple addition often leads to parameter interference and performance degradation. 
TIES-Merging \cite{yadav2023ties} resolves interference by signing and pruning small parameters. 
DARE \cite{yu2023dare} uses random dropping and scaling to merge weights safely. 
Other approaches align neural representations prior to merging, such as Git Re-Basin \cite{ainsworth2022git} and ZipIt! \cite{stoica2023zipit}, or use Fisher Information \cite{matena2021merging} to weight parameter coordinates. 
Despite their elegance, static merging compress distinct adapters into a single frozen state, forcing a permanent compromise in capacity and causing severe accuracy loss under high task diversity.

\subsection{Dynamic Routing and Mixture-of-Experts}
To preserve task-specific specialization, dynamic routing serves multiple adapters by routing input samples on-the-fly. 
Sparsely-Gated Mixture-of-Experts (MoE) \cite{shazeer2017outrageously, fedus2022switch, riquelme2021scaling} route tokens through specialized expert sub-networks using trainable gating functions. 
While effective, standard MoE requires specialized pre-training and substantial parameter overhead, making it unsuitable for post-hoc ensembling of independently fine-tuned adapters.
For adapter-level ensembling, frameworks like S-LoRA \cite{sheng2023slora} and Punica \cite{chen2023punica} batch concurrent heterogeneous requests, but execute each expert in isolation, leading to high memory footprint. 
LoraHub \cite{huang2024lorahub} optimizes a linear mixture of LoRAs on a small validation set, but the resulting mixture is static during inference. 
SABLE and SPS-ZCA address these challenges by enabling training-free, single-pass activation-space blending. 
Specifically, SPS-ZCA uses Zero-Shot Centroid Alignment (ZCA) on intermediate representations to perform hard/temperature-scaled routing. 
However, all these methods represent static feedforward projections where expert activations do not interact. They lack feedback loops, making them highly vulnerable to input noise and domain drift, which blur representations and cause severe misrouting.

\subsection{Connectionist Roots of Lateral Inhibition}
Our competitive exclusion mechanism, which suppresses out-of-domain noise-driven activations, is deeply rooted in the foundational connectionist literature. The concept of lateral inhibition---where active units suppress their neighbors to sharpen representations and isolate a single winner---is a fundamental mechanism in biological brains and early neural network architectures. Key historical connectionist frameworks utilizing such competitive dynamics include Self-Organizing Maps (SOM) \cite{kohonen1982self}, Hopfield Networks \cite{hopfield1982neural}, and Adaptive Resonance Theory (ART) \cite{carpenter1987art}. These classic architectures employ recurrent inhibition or winner-take-all dynamics to perform pattern completion, clustering, and noise filtering. 

We contextualize ESM-LVC as a modern, transformer-era generalization of these principles. Rather than applying lateral inhibition to individual neurons or class logits, we apply non-linear Lotka-Volterra competition-cooperation dynamics to high-dimensional specialized adapter channels (experts) in the activation-space of a frozen backbone during inference. This links classical attractor-network dynamics directly to modern Parameter-Efficient Fine-Tuning serving workloads.

\subsection{Ecological and Dynamical Systems in Deep Learning}
Dynamical systems have a rich history in deep learning, primarily in establishing connections between Residual Networks and Ordinary Differential Equations (ODEs) \cite{he2016deep}, leading to Neural ODEs. 
In biological sciences, the Lotka-Volterra equations represent a cornerstone of mathematical ecology, modeling the population dynamics of species experiencing competition, mutualism, or predation. 
While Lotka-Volterra models have been applied theoretically to study optimization trajectories and neural network stability, they have never been integrated into the activation space of a transformer to dynamically regulate model ensembling. 
By framing expert ensembling as a self-organizing symbiotic ecosystem governed by Lotka-Volterra dynamics, ESM-LVC is the first to bridge mathematical ecology and multi-task model serving, introducing a training-free feedback loop that resolves the limitations of static feedforward routing.
